\section{Métodos}
\label{sec:metodos}

\subsection{NSGA-II adaptado a rutas}
\label{sec:nsga2}

Se ha implementado NSGA-II \citep{deb2002} desde cero, adaptado a secuencias de longitud variable (\code{src/algorithms/genetic\_algorithm.py}). Se eligió porque aproxima todo el frente en una sola ejecución, trata las restricciones sin penalizaciones y es la referencia habitual con la que comparar en optimización multiobjetivo.

\paragraph{Representación.} Un individuo es una secuencia $R = (g_1,\dots,g_k)$ de POIs permitidos ($g_j \in P \setminus F$) sin repeticiones. El hotel no forma parte del cromosoma: se añade al principio y al final al evaluar~\eqref{eq:secuencia}. La longitud varía durante la evolución. En los frentes finales sin preferencias las rutas tienen entre 1 y 7 paradas en Santa Catalina y entre 1 y 6 en Lopesan.

\paragraph{Cruce.} Se elige un punto de corte en el primer progenitor y se copia su prefijo. Después se añaden, en su orden, los genes del segundo progenitor que no estén ya en el prefijo, y el resultado se recorta a 12 genes. El segundo hijo se construye intercambiando los progenitores. El operador conserva subrutas de ambos padres sin repetir POIs. Tiende a producir hijos más largos que los padres; el crecimiento lo frenan el recorte, la mutación de eliminación y, sobre todo, la presión de la restricción de tiempo.

\paragraph{Mutación.} Se aplica uno de cuatro operadores, elegido con estas probabilidades: inserción de un POI permitido no visitado en una posición aleatoria (40 \%), eliminación de una parada si la ruta tiene más de una (30 \%), intercambio de dos paradas (20 \%) e inversión de un segmento (10 \%), un movimiento de tipo 2-opt. Los dos últimos solo cambian el orden, lo que importa porque la matriz es asimétrica.

\paragraph{Restricción de tiempo: reglas de Deb.} Sea $\mathrm{viol}(R) = \max(0, T(R) - \Tmax)$. Según \citet{deb2000}, una solución factible domina a cualquier no factible; entre dos no factibles domina la de menor violación; y entre dos factibles se aplica la dominancia de Pareto:
\begin{equation}
  R_A \succ R_B \iff \bigl(S_A \ge S_B \wedge T_A \le T_B\bigr) \wedge \bigl(S_A > S_B \vee T_A < T_B\bigr).
\end{equation}
Así no hace falta penalizar los objetivos. Aunque el 77 \% de las rutas iniciales supera las 8 horas en Santa Catalina, la población es completamente factible en unas cinco generaciones.

\paragraph{Reparación para las preferencias.} Las restricciones~\eqref{eq:preferencias} se garantizan con la función \code{\_repair\_chromosome}, que en $O(k)$:
\begin{enumerate}
  \item recorre la secuencia y conserva cada gen solo si no está prohibido ni repetido, manteniendo el orden relativo;
  \item si falta la visita obligatoria $m$, la inserta en una posición aleatoria.
\end{enumerate}
Se aplica a los individuos iniciales, tras el cruce y tras la mutación, de modo que todo individuo evaluado cumple $R \subseteq P \setminus F$, no tiene repeticiones y contiene $m$. Por inducción sobre las generaciones, todas las soluciones devueltas cumplen las preferencias. La mutación de eliminación nunca elige $m$, para no destruir su posición optimizada. Además, el problema rechaza al construirse una especificación contradictoria ($m = h$ o $m \in F$). Si la visita obligatoria no cabe en $\Tmax$, todos los individuos son no factibles y el frente devuelto es vacío, de acuerdo con la condición de factibilidad de la sección~\ref{sec:preferencias}.

\paragraph{Pseudocódigo.} El Algoritmo~\ref{alg:nsga2} resume el método.

\begin{algorithm}[H]
\caption{NSGA-II para el problema de orientación biobjetivo}
\label{alg:nsga2}
\small
\begin{algorithmic}[1]
\Require instancia $(V, s, v, t, \Tmax, F, m)$, población $N$, generaciones $G$, probabilidades $p_c$, $p_m$
\Ensure soluciones factibles no dominadas
\State $P_0 \gets$ $N$ secuencias de 3 a 7 POIs permitidos elegidos al azar, reparadas y evaluadas
\For{$g = 1$ \textbf{to} $G$}
  \State $Q_g \gets \emptyset$
  \While{$|Q_g| < N$}
    \State $p_1, p_2 \gets$ torneos binarios en $P_{g-1}$ (menor rango; a igual rango, mayor aglomeración)
    \State $c_1, c_2 \gets$ cruce$(p_1, p_2)$ con probabilidad $p_c$; si no, copias de $p_1, p_2$
    \State reparar $c_1, c_2$; mutar cada uno con probabilidad $p_m$ (la mutación repara)
    \State evaluar $c_1, c_2$ según~\eqref{eq:secuencia}; $Q_g \gets Q_g \cup \{c_1, c_2\}$
  \EndWhile
  \State $R_g \gets P_{g-1} \cup Q_g$ \Comment{$2N$ individuos}
  \State $\mathcal{F}_1, \mathcal{F}_2, \dots \gets$ ordenación no dominada de $R_g$ con las reglas de Deb
  \State $P_g \gets$ frentes completos mientras quepan en $N$; el último, por aglomeración descendente
\EndFor
\State \Return soluciones factibles y distintas de $\mathcal{F}_1$ en $P_G$
\end{algorithmic}
\end{algorithm}

\paragraph{Complejidad.} Con $M = 2$ objetivos, población $N$ y longitud de cromosoma acotada por $k_{\max} = 12$, en cada generación:
\begin{itemize}
  \item crear, reparar y evaluar $N$ hijos cuesta $O(N\,k_{\max})$, porque cada operación recorre como mucho $k_{\max}$ genes;
  \item la ordenación no dominada compara todos los pares de la población combinada de $2N$ individuos: $O(M N^2)$;
  \item la distancia de aglomeración ordena cada frente por cada objetivo: $O(M N \log N)$.
\end{itemize}
El coste por generación es, por tanto, $O(MN^2 + MN\log N + N k_{\max}) = O(MN^2)$, y el de una ejecución completa $O(G M N^2)$. Con los valores usados, una ejecución tarda unos pocos segundos.

\paragraph{Parámetros.} La tabla~\ref{tab:parametros} recoge la configuración. Se fijó de forma empírica y se comprobó después con un estudio de hiperparámetros (sección~\ref{sec:hiperparametros}).

\begin{table}[H]
\centering
\small
\caption{Parámetros de NSGA-II.}
\label{tab:parametros}
\begin{tabularx}{\textwidth}{@{}l c X@{}}
\toprule
Parámetro & Valor & Justificación \\
\midrule
Población $N$ & 100 & Diversidad suficiente para cubrir el frente con un coste $O(MN^2)$ bajo. \\
Generaciones $G$ & 200 & Con la semilla 42 en Santa Catalina, el hipervolumen alcanza el 99 \% de su valor final en la generación 139 y sigue mejorando despacio hasta la 200. \\
Probabilidad de cruce $p_c$ & 0,85 & Valor alto habitual para recombinar subrutas de los padres. \\
Probabilidad de mutación $p_m$ & 0,35 & Alta para explorar un espacio discreto y evitar el estancamiento. \\
Inicialización & 3 a 7 POIs & Rutas aleatorias de tamaño variado; las reglas de Deb corrigen las no factibles. \\
Presupuesto & 20.100 & $N + N \cdot G = 100 + 100 \cdot 200$ evaluaciones. \\
\bottomrule
\end{tabularx}
\end{table}

\subsection{Líneas base}
\label{sec:lineas_base}

Para saber si NSGA-II aporta algo, se compara con cuatro métodos. Los estocásticos consumen exactamente el mismo presupuesto de 20.100 evaluaciones.

\paragraph{Heurística voraz.} Para un presupuesto $B$, parte del hotel y añade en cada paso el POI no visitado $j$ con mayor cociente $s_j / \max(1,\ t_{uj} + v_j)$, siendo $u$ la última parada, siempre que el tiempo acumulado más la vuelta al hotel no supere $B$. Se repite para 29 presupuestos, de 60 a 480 min cada 15 min, y se devuelven las rutas no dominadas. Es determinista y cuesta $O(K_B |P|^2)$ con $K_B = 29$.

\paragraph{Búsqueda aleatoria.} Genera $E = 20.100$ rutas de 2 a 10 POIs elegidos al azar, descarta las que superan $\Tmax$ y devuelve las no dominadas. Su coste asintótico es $O(E\,k_{\mathrm{rs}} + F_{\mathrm{rs}}^2)$, siendo $k_{\mathrm{rs}} \le 10$ el tamaño de muestra y $F_{\mathrm{rs}}$ las soluciones factibles acumuladas antes del filtrado de Pareto.

\paragraph{Suma ponderada.} Para cada peso $w \in \{0{,}1;\ 0{,}2;\ \dots;\ 0{,}9;\ 0{,}95\}$ maximiza, con un algoritmo genético de un solo objetivo,
\begin{equation}
  F_{w}(R) =
  \begin{cases}
    w\, \dfrac{S(R)}{S_{\max}} - (1-w)\, \dfrac{T(R)}{\Tmax} & \text{si } T(R) \le \Tmax,\\[6pt]
    -10^5 - 100\,\bigl(T(R) - \Tmax\bigr) & \text{en otro caso,}
  \end{cases}
\end{equation}
donde $S_{\max} = \sum_{i \in P} s_i$, y se queda con la mejor ruta factible de cada peso.

\paragraph{$\varepsilon$-restricción.} Para 10 cotas equiespaciadas $\varepsilon_k \in [60, 480]$ min resuelve
\begin{equation}
  \max\; S(R) - 10^{-4}\, \frac{T(R)}{\Tmax} \quad \text{sujeto a} \quad T(R) \le \varepsilon_k,
\end{equation}
penalizando las rutas que superan la cota con $-10^5 - 100\,(T(R) - \varepsilon_k)$. El pequeño término en $T$ desempata a favor de la ruta más corta entre las de igual puntuación, para evitar soluciones débilmente eficientes \citep{mavrotas2009}.

\paragraph{Presupuesto de los métodos de escalarización.} Los dos usan el mismo algoritmo genético de un solo objetivo, con los operadores y la reparación de NSGA-II, una población de 30 individuos y 66 generaciones por subproblema. Con 10 subproblemas suman $10 \times (30 + 30 \cdot 66) = 20.100$ evaluaciones. Su coste es $O\bigl(K\,G_{\mathrm{sub}}\,(N_{\mathrm{sub}}\,k_{\max} + N_{\mathrm{sub}} \log N_{\mathrm{sub}})\bigr)$, porque en cada generación se ordena por una aptitud escalar en lugar de por no dominancia. El Algoritmo~\ref{alg:baselines} resume los esquemas de todas las líneas base.

\begin{algorithm}[tbp]
\caption{Líneas base: heurística voraz, búsqueda aleatoria y escalarización GA}
\label{alg:baselines}
\small
\begin{algorithmic}[1]
\Statex \textbf{A) Heurística voraz multiobjetivo} ($B \in \{60, 75, \dots, 480\}$, paso 15 min)
\For{cada presupuesto $B$}
  \State $u \gets h$;\; $t_{\mathrm{acum}} \gets 0$;\; $R \gets ()$;\; $\mathrm{disp} \gets P \setminus F$
  \While{$\exists\, j \in \mathrm{disp} : t_{\mathrm{acum}} + t_{uj} + v_j + t_{jh} \le B$}
    \State $j^* \gets \arg\max_{j \text{ que quepa}} s_j / \max(1,\ t_{uj} + v_j)$;\; añadir $j^*$ a $R$;\; $u \gets j^*$;\; actualizar $t_{\mathrm{acum}}$
  \EndWhile
\EndFor
\State \Return rutas no dominadas entre las generadas
\medskip
\Statex \textbf{B) Búsqueda aleatoria} ($E = 20.100$ evaluaciones)
\State Muestrear $E$ secuencias de 2 a 10 POIs de $P \setminus F$; conservar las factibles ($T \le \Tmax$) y devolver las no dominadas
\medskip
\Statex \textbf{C) Suma ponderada y $\varepsilon$-restricción} ($K = 10$ parámetros $\theta$, $N_{\mathrm{sub}} = 30$, $G_{\mathrm{sub}} = 66$)
\For{cada $\theta \in \{w_1,\dots,w_{10}\}$ o $\{\varepsilon_1,\dots,\varepsilon_{10}\}$}
  \State $\mathit{Pob} \gets$ $N_{\mathrm{sub}}$ secuencias reparadas y evaluadas según~\eqref{eq:secuencia}
  \For{$g = 1$ \textbf{to} $G_{\mathrm{sub}}$}
    \State $Q \gets$ hijos generados por torneo según aptitud($\theta$), cruce, reparación y mutación
    \State $\mathit{Pob} \gets$ $N_{\mathrm{sub}}$ mejores individuos de $\mathit{Pob} \cup Q$ ordenados por aptitud($\theta$)
  \EndFor
  \State guardar la mejor solución factible de $\mathit{Pob}$
\EndFor
\State \Return soluciones no dominadas del conjunto recogido
\end{algorithmic}
\end{algorithm}

\subsection{Rutas recomendadas}
\label{sec:recomendadas}

Un frente con unas 25 rutas no es fácil de leer para un turista. Para ayudar a decidir, de cada frente se extraen tres rutas:
\begin{itemize}
  \item \textbf{Express}: la de mayor puntuación por hora, $S(R)/T(R)$.
  \item \textbf{Equilibrada}: el punto rodilla, es decir, la ruta más cercana al punto ideal (tiempo mínimo y puntuación máxima del frente) tras normalizar ambos objetivos al rango $[0,1]$.
  \item \textbf{Exhaustiva}: la de mayor puntuación dentro de $\Tmax$.
\end{itemize}
